% Notas de aula — Capítulo 20: Previsão Numérica do Tempo
% Curso: Meteorologia Prática (baseado em R. Stull, Practical Meteorology, CC BY-NC-SA 4.0)
% Caixas verdes = conteúdo literal do quadro; linhas verdes = encenação.
\documentclass[11pt]{article}
\usepackage{../comum/preambulo}
\graphicspath{{../figuras/cap20/}}

\title{Capítulo 20 --- Previsão Numérica do Tempo\\
{\large Notas de aula (roteiro de quadro-negro)}}
\author{Curso de Meteorologia Prática}
\date{}

\begin{document}
\maketitle
\thispagestyle{fancy}

\noindent\textbf{Plano sugerido:} 2 aulas de 2\,h.
\emph{Aula 1:} \S\ref{sec:eqs}--\S\ref{sec:param} (as equações, a
grade, esquemas, CFL, parametrizações).
\emph{Aula 2:} \S\ref{sec:caos}--\S\ref{sec:ensemble} (caos e
previsibilidade, assimilação, ensembles).

\section{O sonho de Richardson}\label{sec:eqs}

\begin{noquadro}[abertura]
\textbf{Cap.~20 --- Previsão Numérica (NWP)}\\[2pt]
as 7 equações que o curso inteiro montou:\\
movimento (3) \javimos{10}{Newton girante} $+$ continuidade
\javimos{10}{massa} $+$ energia \javimos{3}{1ª lei} $+$ água
\javimos{4}{balanço de $r_T$} $+$ estado \javimos{1}{gases}\\[2pt]
$=$ um problema de valor inicial: conheça o agora, integre o futuro
\end{noquadro}

Richardson (1922) tentou à mão e errou por 100$\times$ (faltava
filtrar ondas rápidas); Charney, Fjørtoft e Von Neumann acertaram no
ENIAC (1950). Hoje: $\sim$10\,km globais, $10^9$ pontos de grade — e o
mesmo esqueleto conceitual, que este capítulo desmonta em quatro
peças: discretizar, parametrizar, inicializar, quantificar incerteza.

\section{A grade e os esquemas}\label{sec:grade}

\quadro{Discretizar no quadro: $\partial\phi/\partial x \to$
diferenças de vizinhos. Anunciar os dois pecados (difusão e
dispersão) antes de mostrar o Caso~1.}

\begin{formadif}[CFL (T1) e a difusão escondida (T2)]
Von Neumann no upwind ($\phi_j^{n+1} = \phi_j - C(\phi_j -
\phi_{j-1})$, $C = U\Delta t/\Delta x$):
\[ |g|^2 = 1 - 2C(1-C)(1-\cos\theta) \le 1
\iff 0 \le C \le 1\ \ \text{(CFL)}. \]
E a equação modificada revela: o upwind resolve advecção $+$ difusão
com $K_{num} = U\Delta x(1-C)/2$ — as frentes derretem
\javimos{12}{frentes} (N1 mede).
\end{formadif}

\begin{noquadro}[a régua da resolução (T5)]
resolução efetiva $\simeq 7\Delta x$\\
modelo de ``10\,km'' enxerga bem $\gtrsim$70\,km: ciclones sim
\javimos{13}{4000\,km}, cumulus não \javimos{14}{1--10\,km}\\
$\Rightarrow$ o que a grade não vê precisa de \textbf{parametrização}
\end{noquadro}

\section{Parametrizações: a física sub-grade}\label{sec:param}

Tudo que o curso viu ``por dentro'' vira caixa estatística no modelo:
radiação \javimos{2}{Schwarzschild}, microfísica \javimos{7}{gotas e
gelo}, convecção \javimos{14}{cumulus}, turbulência da CLA
\javimos{18}{fluxos turbulentos}, arrasto de ondas de montanha
\javimos{17}{ondas}. Cada uma é um modelo físico reduzido — e a maior
fonte de erro sistemático dos modelos: entender a física original
(Caps.~2--18) é entender \emph{onde} a previsão falha.

\section{Caos: o prazo de validade}\label{sec:caos}

\quadro{Contar a história: Lorenz reroda uma simulação com 3 casas
decimais a menos e o tempo diverge — 1963, o século XX descobre que
determinismo $\neq$ previsibilidade.}

\begin{formadif}[Lorenz e o horizonte (T3; Caso~2)]
\[ \epsilon(t) = \epsilon_0 e^{\lambda t}
\quad\Rightarrow\quad
t_{max} = \frac{1}{\lambda}\ln\frac{E}{\epsilon_0}. \]
Logarítmico em $\epsilon_0$: melhorar o estado inicial 10$\times$
compra só $\sim$2--3 dias. Com erro dobrando a cada $\sim$2 dias, o
teto é $\sim$2 semanas — limite \emph{estrutural}, não computacional.
O tempo é imprevisível; a estatística (clima) não
\veremos{21}{por que o clima é previsível}.
\end{formadif}

\section{Assimilação de dados}\label{sec:assim}

\quadro{Escrever a média ponderada e declarar: ``esta é a fórmula
mais valiosa da meteorologia operacional''.}

\begin{formalivro}[o ganho ótimo (T4; Caso~3)]
\[ x_a = x_b + K(y_o - x_b),\qquad
K^* = \frac{\sigma_b^2}{\sigma_b^2 + \sigma_o^2},\qquad
\frac{1}{\sigma_a^2} = \frac{1}{\sigma_b^2} + \frac{1}{\sigma_o^2}. \]
As precisões \emph{somam}: a análise é melhor que o modelo E que a
observação. Generalizações: 4D-Var, filtro de Kalman por conjunto —
mas o coração é esta média ponderada, ciclada a cada 6\,h sobre
$10^7$ observações \javimos{8}{sondadores} \javimos{9}{rede de
observação}.
\end{formalivro}

O Caso~3 põe um Lorenz caótico na linha com observações ruidosas — e
o N3 encontra o $K$ ótimo empiricamente, batendo com a fórmula.

\section{Ensembles: prever a incerteza}\label{sec:ensemble}

\begin{noquadro}[a mudança filosófica]
1 rodada $=$ ilusão de certeza\\
30--50 rodadas perturbadas $=$ \textbf{distribuição} de futuros\\
dispersão $\approx$ barra de erro; fração de membros $=$
probabilidade\\[2pt]
``70\% de chance de chuva'' É a saída do ensemble (Caso~4; N4 mostra
que o spread prevê o próprio erro)
\end{noquadro}

Verificação fecha o ciclo: índices de acerto, viés, e o ganho de
$\sim$1 dia de previsibilidade por década desde 1980 — a ``revolução
silenciosa'' construída de satélites \javimos{8}{radiâncias},
assimilação e física melhor.

\section{Síntese da aula}

\begin{noquadro}[mapa final --- canto do quadro]
\[ C = \frac{U\Delta t}{\Delta x} \le 1 \qquad
K_{num} = \frac{U\Delta x(1-C)}{2} \qquad
t_{max} = \frac{1}{\lambda}\ln\frac{E}{\epsilon_0} \]
\[ x_a = x_b + K(y_o - x_b) \qquad
\text{ensemble: dispersão} \to \text{probabilidade} \]
discretizar $\cdot$ parametrizar $\cdot$ inicializar $\cdot$
quantificar — o curso inteiro dentro de um computador
\end{noquadro}

\section{Exercícios complementares}\label{sec:exercicios}

Soluções (professor) em \texttt{cap20\_solucoes.ipynb}.

\subsection*{Teóricos}

\begin{enumerate}[label=\textbf{T\arabic*.},itemsep=4pt]
  \item Faça a análise de von Neumann do esquema upwind e deduza a
        condição CFL $0 \le C \le 1$; interprete fisicamente.
  \item Pela equação modificada, mostre que o upwind embute difusão
        $K_{num} = U\Delta x(1-C)/2$ e avalie para um modelo de
        10\,km com $C = 0{,}5$.
  \item Mostre que o horizonte de previsibilidade é
        $t_{max} = \lambda^{-1}\ln(E/\epsilon_0)$ e discuta por que
        melhorar observações rende tão pouco tempo extra.
  \item Minimize a variância da análise e obtenha
        $K^* = \sigma_b^2/(\sigma_b^2+\sigma_o^2)$; mostre que as
        precisões somam.
  \item Explique a ``resolução efetiva $\simeq 7\Delta x$'' pelos
        erros de fase das ondas curtas na grade, e conecte com a
        necessidade de parametrizar convecção.
\end{enumerate}

\subsection*{Numéricos (no notebook)}

\begin{enumerate}[label=\textbf{N\arabic*.},itemsep=4pt]
  \item Meça $K_{num}$ do upwind por ajuste gaussiano e compare com a
        fórmula do T2.
  \item Estime o expoente de Lyapunov do Lorenz com 20 pares de
        gêmeos (média $\pm$ desvio).
  \item Varra o ganho $K$ no ciclo de assimilação e compare o ótimo
        empírico com a fórmula do T4.
  \item Verifique a relação dispersão--erro do ensemble em 50
        previsões: o ensemble sabe quando vai errar?
\end{enumerate}

\vspace{1em}
\noindent\rule{\linewidth}{0.4pt}\\
{\scriptsize Material derivado de R.~Stull, \emph{Practical Meteorology}
(CC BY-NC-SA 4.0); este documento herda a mesma licença.}

\end{document}
